\scanpage{12}
Giving $q$ amounts to : \unsure{the|the} giving of
$\sigma\in
\tikz[baseline=(aut.base)]{
  \node[inner sep=0pt] (aut) {$\operatorname{Aut}^{\pm1}$};
  % Two tiny detached traces above the left of A; writing or surface marks uncertain.
  \node[anchor=south,inner sep=0pt] at ([xshift=-.12em,yshift=.12em]aut.north west)
  {\unsure{\tikz[x=1mm,y=1mm]{\fill (0,.65) circle (.12);\fill[opacity=.45] (0,0) ellipse (.10 and .16);}}};
}(\unsure{G|Z},T)$
\[
\left\{\begin{aligned}p(\sigma)&=-1\\ \sigma^2&=1\end{aligned}\right.
\]
\[
\begin{aligned}
&q_t(-1)=\sigma_t\\[.6ex]
&\sigma_{tg}=g\,\sigma_t\,g^{-1}
=g\,\underbrace{\sigma_t\,g^{-1}\,\sigma_t^{-1}}_{g}\,\sigma_t
=g^2\,\sigma_t=\sigma_t\,g^{-2}\\
&\hspace{15em}\boxed{\sigma_{tg}=g^2\sigma_t=\sigma_t g^{-2}}
\end{aligned}
\]

\uline{Corollary} If $n$ odd, $t\longmapsto\deleted{\unsure{q|g}}\sigma_t$ is a bijection of $T$ with the set \deleted{\unsure{of the}} $p^{-1}\{-1\}$; thus if one has a \unsure{polygon|polygon} with an odd number of sides, the functio\textsuperscript{n} $t\longmapsto\sigma_t$ which, to every vertex of the polygon, assigns the unique orientation-reversing automorphism which fixes this vertex, is bijection.

\uline{NB} If $n$ even, \unsure{\deleted{it}|it} \unsure{this} is not bijective, \deleted{\unsure{but}} \deleted{\unsure{for all}} its fibers are empty or of cardinality 2\ldots

% slanted left margin note; the line nodes preserve the cramped oblique geometry.
\reversemarginpar
\margin{\vspace*{-2.7\baselineskip}\hspace*{-1.05cm}%
\begin{tikzpicture}[x=.007cm,y=.007cm,baseline=(0,0)]
  \node[anchor=base west,rotate=34,inner sep=0pt] at (31,-115) {$\{u,\bar u\}\subset Z$};
  \node[anchor=base west,rotate=34,inner sep=0pt] at (15,-183) {opposite generators of order};
  \node[anchor=base west,rotate=34,inner sep=0pt] at (10,-239) {$n=2m$, then};
  \node[anchor=base west,rotate=34,inner sep=0pt] at (30,-291) {$u^m=\bar u^m=s$, \uline{unique}};
  \node[anchor=base west,rotate=34,inner sep=0pt] at (23,-348) {\uline{\unsure{element|element} of order} 2 of $Z$};
  \node[anchor=base west,rotate=34,inner sep=0pt] at (25,-412) {called \unsure{symmetry|symmetry}};
  \node[anchor=base west,rotate=34,inner sep=0pt] at (25,-465) {My $\sigma_t=\sigma_{t'}$};
  \node[anchor=base west,rotate=34,inner sep=0pt] at (35,-520) {\unsure{if and only if|sss} $t,t'$ are \unsure{conjugate}};
  \node[anchor=base west,rotate=34,inner sep=0pt] at (43,-565) {modulo $\{1,s\}$};
  % The separate upright and the thinner curved boundary are not mathematical arrows.
  \draw[line width=.45pt] (313,-5) -- (319,-179);
  \draw[line width=.28pt] (285,-73) -- (307,-154)
    .. controls (337,-187) and (393,-193) .. (427,-193)
    .. controls (377,-221) and (337,-255) .. (335,-281)
    .. controls (342,-359) and (375,-401) .. (334,-476)
    .. controls (291,-536) and (222,-603) .. (152,-665);
\end{tikzpicture}}
\normalmarginpar

\uline{Opposite} \uline{vertex}, canonical symmetry

\uline{Parenthesis}\qquad Fiber product of sets

\diagram{\begin{tikzpicture}
\node[anchor=north west,inner sep=0pt] (square) at (0,0) {\begin{tikzcd}[ampersand replacement=\&,column sep=2em,row sep=2.4em]
E\arrow[r,"p"] \& \unsure{S|B}\\
P\arrow[u,dashed,"q'"]\arrow[r,dashed,"p'"] \& F\arrow[u,"q"']
\end{tikzcd}};
% The set description begins at the level of the lower row of the square.
\node[anchor=north west,inner sep=0pt] at (4.5,-1.5) {$\begin{aligned}
&P\subset E\times F\\[-.4em]
&\hspace{.3em}\rotatebox{90}{$=$}\\[-.4em]
&\{(x,y)\mid px=qy\}
\end{aligned}$};
\end{tikzpicture}}

% Detached pair of points to the far right of the set display; no punctuation function assigned.
\margin{\raisebox{1.2em}{\unsure{\tikz[x=1mm,y=1mm]{\fill (.45,.8) circle (.12);\fill (0,0) circle (.12);}}}}

% Detached, definite ink stroke above Solution; no lexical value is assigned.
\figure{\makebox[\linewidth][l]{\begin{tikzpicture}[x=1cm,y=1cm]
\fill (0.025,0.29) .. controls (-0.015,0.28) and (-0.012,0.19) .. (0.018,0.15)
 .. controls (0.038,0.12) and (0.039,0.075) .. (0.012,0.055)
 .. controls (-0.005,0.039) and (-0.036,0.055) .. (-0.043,0.037)
 .. controls (-0.044,0.008) and (0.003,-0.004) .. (0.028,0.022)
 .. controls (0.071,0.06) and (0.06,0.12) .. (0.05,0.17)
 .. controls (0.072,0.23) and (0.062,0.27) .. (0.025,0.29) -- cycle;
\end{tikzpicture}}}

Solution of a universal problem in the category of sets :

If $E,F,S$ \unsure{groups|groups}, $p,q$ \unsure{group|group} homomorphisms, then $P$ a \unsure{group|group} (and solution of the analogous universal problem \unsure{in|in} category of \unsure{groups|groups})
\[
\operatorname{Ker}p\sim\operatorname{Ker}p',\quad
\operatorname{Ker}q\sim\operatorname{Ker}q'
\]

If $p$\tikz[baseline=(surj.base)]{
  \node[anchor=base west,inner sep=0pt] (surj) at (0,0) {surjective};
  \node[anchor=base west,inner sep=0pt] at (-.05, .39) {\inserted{$:E\mathbin{\text{---}}S$}};
  \draw[line width=.3pt] (.02,.12) .. controls (.24,.32) and (.65,.30) .. (.94,.40);
}\ bijective) then $P\mathbin{\text{---}}F$ surjective bijective).

[Language of Cartesian
\tikz[baseline=(old.base)]{
  \node[anchor=base,inner sep=0pt] (old) at (0,0) {\deleted{\unsure{diagrams}}};
  \node[anchor=base,inner sep=0pt] at (.12,.38) {\inserted{squares}};
}
]

Inverse image of ex. of \unsure{groups|group|group|groups}
% archive sheet number: 11
